\clearpage
\renewcommand{\sectioncolor}{slate}
\section{Verification ledger --- what was checked, and how}
\markboth{Verification ledger}{}

Every claim of fact in this tutorial was checked on your machine on \textbf{19 September 2026}.
Where a check could not be made, it says so.

\begin{center}\footnotesize
\setlength{\tabcolsep}{5pt}\renewcommand{\arraystretch}{1.22}
\begin{tabular}{@{}>{\raggedright\arraybackslash}p{5.0cm} >{\raggedright\arraybackslash}p{5.6cm} >{\raggedright\arraybackslash}p{5.0cm}@{}}
\toprule
\rowcolor{slate!13}
\textbf{claim} & \textbf{how checked} & \textbf{result}\\
\midrule
Notebook has 22 cells, 10 of them code & \bk{json.load} over the \bk{.ipynb} &
 \textcolor{bio}{\textbf{22 cells, 10 code, 12 markdown}}\\
\rowcolor{slate!5}
Notebook runs clean under Sage & scan every cell's outputs for \bk{output\_type=='error'} &
 \textcolor{bio}{\textbf{0 errors}}, 10 cells executed\\
Kernel is Sage 10.9 & \bk{metadata.kernelspec} &
 \bk{SageMath 10.9} / \bk{sagemath-10.9}\\
\rowcolor{slate!5}
Eight proof books present & \bk{find} over \bk{Maths\_Proof\_Skill/} &
 8 PDFs, 6 folders, \textbf{22 MB}\\
Copies are byte-identical to sources & MD5 of every copy vs.\ its original &
 \textcolor{bio}{\textbf{8/8 match}}; originals untouched\\
\rowcolor{slate!5}
Quotations in §2--§4 & read from \emph{your} PDFs, not from memory &
 Devlin §3.1; D\&L Def.~2.1.15 and §2.3.6; Thurston; Tao\\
Cell numbers cited in §5--§8 & re-extracted from the notebook and cross-checked &
 corrected once during drafting --- see note below\\
\bottomrule
\end{tabular}
\end{center}

\begin{warnbox}[title={One correction made while writing this document}]
An early draft cited the CAS-failure demonstration as cell~19. Re-extracting the cell list showed
cell~19 is the \emph{theorem verification} and the CAS demonstration is \textbf{cell~21}. Every
\textcolor{proof}{\textbf{CELL~$n$}} tag in this document was then re-derived from the notebook file
rather than from the draft. \textbf{This is the same discipline the document argues for}: a claim
that was only plausible got checked against the source, and it was wrong.
\end{warnbox}

\subsection{Three things this tutorial deliberately does not claim}

\begin{enumerate}\setlength{\itemsep}{4pt}
\item \textbf{That you can now write proofs.} You can now \emph{read} six techniques and recognise
 which one a proof is using. Writing them is weeks 2--4, and it is exercises, not reading.
\item \textbf{That the notebook's Sage checks prove anything.} They do not, and the notebook says so
 in capitals at cell~17. Every proof in it is in the markdown.
\item \textbf{That the four-week route is enough for Axler.} It gets you to the point where Axler
 §7.B is readable. \emph{Readable}, not easy --- that is Week 1 of the twelve-week partition-function
 plan, and a different document.
\end{enumerate}

\vfill
\begin{center}
\begin{tcolorbox}[enhanced,width=0.94\textwidth,colback=proof!6,colframe=proof,
 boxrule=1.2pt,arc=5pt,left=12pt,right=12pt,top=10pt,bottom=10pt]
\centering
{\large\bfseries\color{proof} The one sentence to keep}\\[7pt]
\begin{minipage}{0.88\textwidth}\centering
Computation tells you \emph{that} something is true for the cases you ran.\\
A proof tells you \emph{why} it is true, for all the cases you will never run.\\[6pt]
{\footnotesize\color{slate}Devlin: \emph{``a proof of a statement must explain why that statement is
true.''}\\
Thurston: \emph{``a continuing desire for human understanding of a proof, in addition to knowledge
that the theorem is true.''}}
\end{minipage}
\end{tcolorbox}
\end{center}
\vfill
